% GENERATED FILE. Edit canonical lessons, then run npm run generate:books.
% canonical-source-hash: fnv1a64:4d12bfc149346b28
% canonical-lessons: TA-C464-capalam, TA-C464-virakti, TA-C464-tontaravu, TA-C464-manratu, TA-C464-alu

\chapter{Temptation and Bother}
\label{ch:ta-temptation-and-bother}
\hlchaptermodality{464}

\begin{chapteropening}
\textbf{By the end of this chapter:} \emph{I can talk about temptation, despair, pleading and being fed up.}

Complete the last lesson of chapter 464: I can talk about temptation, despair, pleading and being fed up.
\end{chapteropening}

\section[\texorpdfstring{\ta{சபலம்} (capalam)}{சபலம் (capalam)}]{\ta{சபலம்} (capalam) — temptation; a weak urge}

\label{lesson:TA-C464-capalam}



\begin{quote}
Before the new one: say the Tamil for emptiness, then the Tamil for condolence.
\end{quote}

\subsection*{You'll want to know: \ta{சபலம்}}
\textbf{\ta{சபலம்}} — \emph{capalam} — \textquotedblleft{}temptation; a weak urge\textquotedblright{}.

\begin{quote}
\ta{இனிப்பைப்} \ta{பார்த்ததும்} \ta{எனக்குச்} \ta{சபலம்} \ta{வந்தது}.
\end{quote}

\begin{quote}
\emph{iṉippaip pārttatum eṉakkuc capalam vantatu.}
\end{quote}

\begin{quote}
When I saw the sweets, I was tempted.
\end{quote}

Dative experiencer: \ta{எனக்குச்} \ta{சபலம்} \ta{வந்தது}, temptation came to me.

\begin{cousinweb}
\textbf{\ta{சபலம்}} (\emph{capalam}) is borrowed from Sanskrit \emph{capala}, \textquotedblleft{}trembling, fickle\textquotedblright{} (CDIAL 4672): in Tamil, a wavering will.
\end{cousinweb}

\begin{grammarlens}[title={What you've built}]
One more word for describing what you lived and felt.
\end{grammarlens}

\subsection*{Your turn}
\begin{itemize}
\raggedright
  \item \emph{Say it:} \emph{capalam}
  \item \emph{Say it:} \emph{capalam}, once more
  \item \emph{Say it:} \emph{iṉippaip pārttatum eṉakkuc capalam vantatu.} — When I saw the sweets, I was tempted.
  \item \emph{Recall:} say the Tamil for emptiness, then the Tamil for condolence, then say \emph{capalam} again
\end{itemize}

\begin{tcolorbox}[breakable,colback=teal!4,colframe=teal!35!black,title={Before you move on}]
What does \ta{சபலம்} mean? (\textquotedblleft{}Temptation; a weak urge\textquotedblright{}.) Say it once more.
\end{tcolorbox}
\section[\texorpdfstring{\ta{விரக்தி} (virakti)}{விரக்தி (virakti)}]{\ta{விரக்தி} (virakti) — despair; disillusionment}

\label{lesson:TA-C464-virakti}



\begin{quote}
Before the new one: say the Tamil for condolence, then the Tamil for temptation.
\end{quote}

\subsection*{You'll want to know: \ta{விரக்தி}}
\textbf{\ta{விரக்தி}} — \emph{virakti} — \textquotedblleft{}despair; disillusionment\textquotedblright{}.

\begin{quote}
\ta{பல} \ta{முறை} \ta{தோற்றதால்} \ta{அவருக்கு} \ta{விரக்தி} \ta{வந்தது}.
\end{quote}

\begin{quote}
\emph{pala muṟai tōṟṟatāl avarukku virakti vantatu.}
\end{quote}

\begin{quote}
Having failed many times, he fell into despair.
\end{quote}

\ta{விரக்தியான} \ta{பேச்சு}, bitter, hopeless talk.

\begin{cousinweb}
\textbf{\ta{விரக்தி}} (\emph{virakti}) is borrowed from Sanskrit \emph{virakti}, \textquotedblleft{}change of feeling, loss of interest\textquotedblright{} (Monier-Williams).
\end{cousinweb}

\begin{grammarlens}[title={What you've built}]
One more word for describing what you lived and felt.
\end{grammarlens}

\subsection*{Your turn}
\begin{itemize}
\raggedright
  \item \emph{Say it:} \emph{virakti}
  \item \emph{Say it:} \emph{virakti}, once more
  \item \emph{Say it:} \emph{pala muṟai tōṟṟatāl avarukku virakti vantatu.} — Having failed many times, he fell into despair.
  \item \emph{Recall:} say the Tamil for condolence, then the Tamil for temptation, then say \emph{virakti} again
\end{itemize}

\begin{tcolorbox}[breakable,colback=teal!4,colframe=teal!35!black,title={Before you move on}]
What does \ta{விரக்தி} mean? (\textquotedblleft{}Despair; disillusionment\textquotedblright{}.) Say it once more.
\end{tcolorbox}
\section[\texorpdfstring{\ta{தொந்தரவு} (tontaravu)}{தொந்தரவு (tontaravu)}]{\ta{தொந்தரவு} (tontaravu) — a nuisance; bother}

\label{lesson:TA-C464-tontaravu}



\begin{quote}
Before the new one: say the Tamil for temptation, then the Tamil for despair.
\end{quote}

\subsection*{You'll want to know: \ta{தொந்தரவு}}
\textbf{\ta{தொந்தரவு}} — \emph{tontaravu} — \textquotedblleft{}a nuisance; bother\textquotedblright{}.

\begin{quote}
\ta{பக்கத்து} \ta{வீட்டு} \ta{நாய்} \ta{இரவில்} \ta{எங்களுக்குத்} \ta{தொந்தரவு} \ta{செய்தது}.
\end{quote}

\begin{quote}
\emph{pakkattu vīṭṭu nāy iravil eṅkaḷukkut tontaravu ceytatu.}
\end{quote}

\begin{quote}
The neighbour’s dog bothered us at night.
\end{quote}

\ta{தொந்தரவு} \ta{செய்}, to bother someone, who takes -\ta{க்கு}.

\begin{cousinweb}
\textbf{\ta{தொந்தரவு}} (\emph{tontaravu}) is native Dravidian (DEDR 3506); the entry adds the verb \ta{தொந்தரி} \emph{tontari}, \textquotedblleft{}to vex\textquotedblright{}.
\end{cousinweb}

\begin{grammarlens}[title={What you've built}]
One more word for describing what you lived and felt.
\end{grammarlens}

\subsection*{Your turn}
\begin{itemize}
\raggedright
  \item \emph{Say it:} \emph{tontaravu}
  \item \emph{Say it:} \emph{tontaravu}, once more
  \item \emph{Say it:} \emph{pakkattu vīṭṭu nāy iravil eṅkaḷukkut tontaravu ceytatu.} — The neighbour’s dog bothered us at night.
  \item \emph{Recall:} say the Tamil for temptation, then the Tamil for despair, then say \emph{tontaravu} again
\end{itemize}

\begin{tcolorbox}[breakable,colback=teal!4,colframe=teal!35!black,title={Before you move on}]
What does \ta{தொந்தரவு} mean? (\textquotedblleft{}A nuisance; bother\textquotedblright{}.) Say it once more.
\end{tcolorbox}
\section[\texorpdfstring{\ta{மன்றாடு} (maṉṟāṭu)}{மன்றாடு (maṉṟāṭu)}]{\ta{மன்றாடு} (maṉṟāṭu) — to plead; to beg earnestly}

\label{lesson:TA-C464-manratu}



\begin{quote}
Before the new one: say the Tamil for despair, then the Tamil for a nuisance.
\end{quote}

\subsection*{You'll want to know: \ta{மன்றாடு}}
\textbf{\ta{மன்றாடு}} — \emph{maṉṟāṭu} — \textquotedblleft{}to plead; to beg earnestly\textquotedblright{}.

\begin{quote}
\ta{வேலை} \ta{வேண்டும்} \ta{என்று} \ta{அவர்} \ta{முதலாளியிடம்} \ta{மன்றாடினார்}.
\end{quote}

\begin{quote}
\emph{vēlai vēṇṭum eṉṟu avar mutalāḷiyiṭam maṉṟāṭiṉār.}
\end{quote}

\begin{quote}
He pleaded with his employer for work.
\end{quote}

Past (I …): \ta{மன்றாடினேன்}, \emph{maṉṟāṭiṉēṉ}.

The person pleaded with takes -\ta{இடம்}.

\begin{cousinweb}
The origin of \textbf{\ta{மன்றாடு}} (\emph{maṉṟāṭu}) is not settled; it ends like verbs built with \ta{ஆடு}, but the first part is unclear.
\end{cousinweb}

\begin{grammarlens}[title={What you've built}]
One more word for describing what you lived and felt.
\end{grammarlens}

\subsection*{Your turn}
\begin{itemize}
\raggedright
  \item \emph{Say it:} \emph{maṉṟāṭu}
  \item \emph{Say it:} \emph{maṉṟāṭu}, once more
  \item \emph{Say it:} \emph{vēlai vēṇṭum eṉṟu avar mutalāḷiyiṭam maṉṟāṭiṉār.} — He pleaded with his employer for work.
  \item \emph{Recall:} say the Tamil for despair, then the Tamil for a nuisance, then say \emph{maṉṟāṭu} again
\end{itemize}

\begin{tcolorbox}[breakable,colback=teal!4,colframe=teal!35!black,title={Before you move on}]
What does \ta{மன்றாடு} mean? (\textquotedblleft{}To plead; to beg earnestly\textquotedblright{}.) Say it once more.
\end{tcolorbox}
\section[\texorpdfstring{\ta{அலு} (alu)}{அலு (alu)}]{\ta{அலு} (alu) — to be fed up; to grow tired of}

\label{lesson:TA-C464-alu}



\begin{quote}
Before the new one: say the Tamil for a nuisance, then the Tamil for to plead.
\end{quote}

\subsection*{You'll want to know: \ta{அலு}}
\textbf{\ta{அலு}} — \emph{alu} — \textquotedblleft{}to be fed up; to grow tired of\textquotedblright{}.

\begin{quote}
\ta{தினமும்} \ta{அதே} \ta{உணவைச்} \ta{சாப்பிட்டு} \ta{எனக்கு} \ta{அலுத்துவிட்டது}.
\end{quote}

\begin{quote}
\emph{tiṉamum atē uṇavaic cāppiṭṭu eṉakku aluttuviṭṭatu.}
\end{quote}

\begin{quote}
I got fed up eating the same food every day.
\end{quote}

Past (it …): \ta{அலுத்தது}, \emph{aluttatu}.

Dative experiencer: \ta{எனக்கு} \ta{அலுத்தது}, I am tired of it. \ta{அலுப்பு}, weariness.

\begin{cousinweb}
\textbf{\ta{அலு}} (\emph{alu}) is native Dravidian (DEDR 236), to be worn out by work or care; the same entry gives \ta{அலுப்பு} \emph{aluppu}, \textquotedblleft{}weariness\textquotedblright{}.
\end{cousinweb}

\begin{grammarlens}[title={What you've built}]
One more word for describing what you lived and felt.
\end{grammarlens}

\subsection*{Your turn}
\begin{itemize}
\raggedright
  \item \emph{Say it:} \emph{alu}
  \item \emph{Say it:} \emph{alu}, once more
  \item \emph{Say it:} \emph{tiṉamum atē uṇavaic cāppiṭṭu eṉakku aluttuviṭṭatu.} — I got fed up eating the same food every day.
  \item \emph{Recall:} say the Tamil for a nuisance, then the Tamil for to plead, then say \emph{alu} again
\end{itemize}

\begin{tcolorbox}[breakable,colback=teal!4,colframe=teal!35!black,title={Before you move on}]
What does \ta{அலு} mean? (\textquotedblleft{}To be fed up; to grow tired of\textquotedblright{}.) Say it once more.
\end{tcolorbox}
